Question Five (15 Marks)
(A) Compare and contrast Emitter bias and voltage divider bias configurations in a BJT transistor. [5 Marks]
While both configurations are used to establish a stable DC operating point (Q-point) for a BJT, they achieve stability differently:
Voltage Divider Bias is highly stable and independent of β (beta). It uses a resistive divider network at the base to fix a constant base voltage $V_B$, making the operating point independent of temperature fluctuations or transistor swaps.
Emitter Bias provides moderate stability but requires dual-polarity power supplies ($+V_{CC}$ and $-V_{EE}$) to achieve good stability. It is more dependent on the transistor's inherent parameter variations (β) compared to the voltage divider method.
Feature Voltage Divider Bias Emitter Bias
Supply Voltage Requires only a single DC power supply ($V_{CC}$). Typically requires split/dual DC power supplies ($+V_{CC}$ and $-V_{EE}$).
Q-Point Stability Highest stability; virtually independent of β. Moderate stability; still somewhat reliant on β value matching.
Circuit Complexity Uses more resistors at the base network interface. Simpler base connection, but more complex power routing.
(B) Briefly describe both biasing methods and explain the circuit diagrams and components involved. [4 Marks]
1. Voltage Divider Bias
Description: This setup features a voltage divider string (R₁ and R₂) connected across the single power supply $V_{CC}$ to provide a stable, constant voltage to the BJT base terminal.
Key Components & Layout:
R₁ (pull-up) and R₂ (pull-down) form a voltage divider at the base terminal.
$R_C$ is connected between the collector and $V_{CC}$ to establish the output voltage loop.
$R_E$ is tied between the emitter and ground to provide negative feedback for thermal stabilization.
2. Emitter Bias
Description: This configuration relies on a negative power supply connected directly to the emitter resistor to push a constant current through the transistor channel.
Key Components & Layout:
$R_B$ connects the base directly to the circuit ground reference.
$R_E$ connects the emitter terminal down to a negative supply rail ($-V_{EE}$).
$R_C$ connects the collector up to the positive supply rail ($+V_{CC}$).
(C) Given a voltage divider consisting of resistors R₁ and R₂ connected to a DC supply $V_{CC}$ and a base resistor connecting to the transistor base: [6 Marks]
i. Explain step-by-step how to simplify the voltage divider bias network into its Thévenin equivalent. [2 Marks]
Isolate the Base Terminal: Disconnect the base terminal of the transistor from the junction connecting R₁ and R₂.
Find Thévenin Voltage ($V_{TH}$): Calculate the open-circuit voltage at that node. This is the voltage dropping across resistor R₂ via the voltage divider rule.
Find Thévenin Resistance ($R_{TH}$): Deactivate the DC power source ($V_{CC}$) by replacing it with a short-circuit link to ground. Look back into the base terminal node to determine the parallel equivalent resistance of R₁ and R₂.
ii. Derive expressions for the Thévenin voltage $V_{TH}$ and Thévenin resistance $R_{TH}$. [2 Marks]
Thévenin Voltage equation:
$$V_{TH} = V_{CC} \times \left( \frac{R_2}{R_1 + R_2} \right)$$ 
Thévenin Resistance equation:
$$R_{TH} = R_1 \parallel R_2 = \frac{R_1 \times R_2}{R_1 + R_2}$$ 
iii. Show how this simplification helps in calculating the base current $I_B$, collector current $I_C$, and emitter current $I_E$. [2 Marks]
Once simplified, the complex input network reduces to a simple single-loop loop consisting of $V_{TH}$, $R_{TH}$, the base-emitter junction ($V_{BE}$), and the emitter resistor ($R_E$).
Apply Kirchhoff's Voltage Law (KVL) to the base-emitter input loop:
$$V_{TH} - I_B R_{TH} - V_{BE} - I_E R_E = 0$$ 
Substitute $I_E = (\beta + 1)I_B$ into the equation to calculate Base Current ($I_B$):
$$I_B = \frac{V_{TH} - V_{BE}}{R_{TH} + (\beta + 1)R_E}$$
Use the computed $I_B$ value alongside standard BJT current relationships to easily solve for Collector Current ($I_C$) and Emitter Current ($I_E$):
$$I_C = \beta \cdot I_B$$ 
$$I_E = (\beta + 1) \cdot I_B$$ 
